\chapter{L3 -- Sistemi lineari: matrice di transizione, evoluzione forzata, rappresentazioni esplicite e implicite}

\section{Forma lineare delle mappe di stato e di uscita}

Nella lezione precedente è emersa la proprietà di decomposizione dei sistemi lineari:
\[
\phi(t,t_0,x_0,u)=\phi_\ell(t,t_0,x_0)+\phi_f(t,t_0,u),
\]
\[
g(t,t_0,x_0,u)=g_\ell(t,t_0,x_0)+g_f(t,t_0,u).
\]

Quando il sistema è lineare e gli spazi
\[
X=\mathbb{R}^n,\qquad U=\mathbb{R}^m,\qquad Y=\mathbb{R}^{\ell}
\]
sono finito-dimensionali, la parte \emph{libera} e la parte \emph{istantanea di uscita} si possono rappresentare mediante matrici.

\subsection{Trasformazione d'uscita lineare}

La parte libera dell'uscita dipende linearmente dallo stato:
\[
g_\ell(x(t),t)=C(t)x(t),
\qquad
C(t)\in\mathbb{R}^{\ell\times n}.
\]

La parte forzata dell'uscita dipende linearmente dall'ingresso:
\[
g_f(u(t),t)=D(t)u(t),
\qquad
D(t)\in\mathbb{R}^{\ell\times m}.
\]

Quindi, in forma locale,
\begin{equation}
y(t)=C(t)x(t)+D(t)u(t).
\label{eq:L3_output_local}
\end{equation}

\paragraph{Interpretazione delle dimensioni.}
\begin{itemize}
    \item $x(t)\in \mathbb{R}^n$ è il vettore di stato;
    \item $u(t)\in \mathbb{R}^m$ è il vettore di ingresso;
    \item $y(t)\in \mathbb{R}^{\ell}$ è il vettore di uscita;
    \item $C(t)$ trasforma lo stato in uscita;
    \item $D(t)$ descrive l'eventuale passaggio istantaneo ingresso--uscita.
\end{itemize}

\subsection{Parte libera dello stato e matrice di transizione}

Poiché $\phi_\ell$ è una mappa lineare da $X$ in $X$, si può scrivere
\[
\phi_\ell(t,t_0,x(t_0))=\Phi(t,t_0)x(t_0),
\]
dove
\[
\Phi(t,t_0)\in\mathbb{R}^{n\times n}
\]
è la \emph{matrice di transizione di stato}.

\paragraph{Significato.}
La matrice $\Phi(t,t_0)$ trasforma lo stato iniziale $x(t_0)$ nello stato che il sistema avrebbe al tempo $t$ in assenza di ingresso.

\paragraph{Caso tempo-invariante.}
Se il sistema è lineare tempo-invariante, allora la matrice di transizione dipende solo dal tempo trascorso:
\[
\Phi(t,t_0)=\Phi(t-t_0).
\]

Questa dipendenza solo dalla differenza $t-t_0$ esprime precisamente il fatto che, in un sistema tempo-invariante, non conta \emph{quando} inizio l'evoluzione, ma solo \emph{quanto tempo} è trascorso dall'istante iniziale.

\section{Rappresentazioni esplicite di stato nei sistemi lineari}

\subsection{Tempo continuo}

Nel caso continuo, la rappresentazione esplicita dello stato assume la forma
\begin{equation}
x(t)=\Phi(t,t_0)x(t_0)+\int_{t_0}^{t} H(t,\tau)u(\tau)\,d\tau,
\label{eq:L3_explicit_TC}
\end{equation}
dove
\[
H(t,\tau)\in\mathbb{R}^{n\times m}
\]
è un opportuno \emph{nucleo di evoluzione forzata}.

L'uscita si scrive invece come
\begin{equation}
y(t)=C(t)x(t)+D(t)u(t).
\label{eq:L3_TC_output}
\end{equation}

\paragraph{Lettura della formula.}
La \eqref{eq:L3_explicit_TC} si divide in due contributi:
\begin{itemize}
    \item \emph{risposta libera}
    \[
    \Phi(t,t_0)x(t_0),
    \]
    dovuta solo allo stato iniziale;
    \item \emph{risposta forzata}
    \[
    \int_{t_0}^{t} H(t,\tau)u(\tau)\,d\tau,
    \]
    dovuta solo all'ingresso.
\end{itemize}

\subsection{Tempo discreto}

Nel caso discreto, l'integrale viene sostituito da una somma:
\begin{equation}
x(K)=\Phi(K,K_0)x(K_0)+\sum_{j=K_0}^{K-1} H(K,j)\,u(j),
\label{eq:L3_explicit_TD}
\end{equation}
con
\[
H(K,j)\in\mathbb{R}^{n\times m}.
\]

L'uscita si scrive
\begin{equation}
y(K)=C(K)x(K)+D(K)u(K).
\label{eq:L3_TD_output}
\end{equation}

\paragraph{Interpretazione.}
Per calcolare lo stato al tempo $K$:
\begin{itemize}
    \item porto avanti lo stato iniziale tramite $\Phi(K,K_0)$;
    \item sommo poi tutti i contributi degli ingressi applicati agli istanti
    \[
    j=K_0,K_0+1,\dots,K-1.
    \]
\end{itemize}

\section{Perché nel tempo discreto compare la somma}

Negli appunti questa parte è centrale: si vuole capire da dove nasce in modo naturale la formula
\[
x(K)=\Phi(K,K_0)x(K_0)+\sum_{j=K_0}^{K-1}H(K,j)u(j).
\]

L'idea consiste nel decomporre un ingresso discreto generico in una somma di \emph{ingressi elementari}.

\subsection{Un ingresso discreto su un intervallo finito}

Consideriamo l'ingresso ristretto all'intervallo
\[
[K_0,K)\cap \mathbb{Z}.
\]
Lo indichiamo con
\[
u_{[K_0,K)}(\cdot):T\to U.
\]

Esso è completamente individuato dai valori
\[
u(K_0),\ u(K_0+1),\ \dots,\ u(K-1).
\]

\paragraph{Osservazione importante.}
Anche se $u(\cdot)$ appartiene a uno spazio di funzioni, quando lo osservo solo su un intervallo discreto finito esso viene descritto da un numero finito di campioni. È proprio questo che permette di ricondurre l'evoluzione forzata a una somma finita di contributi elementari.

\begin{figure}[H]
\centering
\begin{tikzpicture}[scale=0.95]
    \draw[->] (0,0) -- (10.2,0) node[right] {$k$};
    \draw[->] (0,0) -- (0,4.2) node[above] {$U$};

    \foreach \x in {1,2,3,4,5,6,7,8,9} {
        \draw (\x,0.12) -- (\x,-0.12);
    }

    \draw[thick,red] (2,0)--(2,1.0);
    \fill[red] (2,1.0) circle (2.3pt);
    \node[red] at (2,-0.45) {$K_0$};
    \node[red] at (0.9,1.15) {$u(K_0)$};

    \draw[thick,green!60!black] (4,0)--(4,2.5);
    \fill[green!60!black] (4,2.5) circle (2.3pt);
    \node[green!60!black] at (4,-0.45) {$K_0+1$};
    \node[green!60!black] at (2.95,2.7) {$u(K_0+1)$};

    \draw[thick,violet] (5,0)--(5,1.8);
    \fill[violet] (5,1.8) circle (2.3pt);
    \node[violet] at (5,-0.45) {$K_0+2$};
    \node[violet] at (3.9,1.95) {$u(K_0+2)$};

    \draw[thick,orange!85!black] (9,0)--(9,3.3);
    \fill[orange!85!black] (9,3.3) circle (2.3pt);
    \node[orange!85!black] at (9,-0.45) {$K-1$};
    \node[orange!85!black] at (7.8,3.5) {$u(K-1)$};

    \fill (6,1.4) circle (2pt);
    \fill (7,3.0) circle (2pt);
    \fill (8,2.4) circle (2pt);
\end{tikzpicture}
\caption{Un ingresso discreto generico è completamente determinato dai suoi valori negli istanti $K_0,\dots,K-1$.}
\end{figure}

\subsection{Decomposizione in ingressi elementari}

Per ogni indice $j\in\{K_0,\dots,K-1\}$ definiamo il segnale elementare
\[
u_j(k)=
\begin{cases}
u(j), & k=j,\\[4pt]
0, & k\neq j.
\end{cases}
\]

Allora si ha la decomposizione
\begin{equation}
u_{[K_0,K)}(\cdot)=\sum_{j=K_0}^{K-1}u_j(\cdot).
\label{eq:L3_input_decomposition}
\end{equation}

\paragraph{Perché funziona.}
Per ogni istante $k$:
\begin{itemize}
    \item solo uno dei segnali $u_j$ può essere diverso da zero, precisamente quello con $j=k$;
    \item quindi la somma dei segnali elementari restituisce esattamente il valore originario $u(k)$.
\end{itemize}

\begin{figure}[H]
\centering
\begin{tikzpicture}[scale=0.9]
    % asse 1
    \begin{scope}[yshift=3.8cm]
        \draw[->] (0,0) -- (7.2,0) node[right] {$k$};
        \draw[->] (0,0) -- (0,2.4) node[above] {$U$};
        \foreach \x in {1,2,3,4,5,6} {\draw (\x,0.12)--(\x,-0.12);}
        \draw[thick,red] (2,0)--(2,1.2);
        \fill[red] (2,1.2) circle (2.1pt);
        \node[red] at (2,-0.35) {$K_0$};
        \node[red] at (0.9,1.45) {$u_{K_0}$};
        \foreach \x in {1,3,4,5,6} {\fill[red] (\x,0) circle (2pt);}
    \end{scope}

    % asse 2
    \begin{scope}[yshift=1.2cm]
        \draw[->] (0,0) -- (7.2,0) node[right] {$k$};
        \draw[->] (0,0) -- (0,2.4) node[above] {$U$};
        \foreach \x in {1,2,3,4,5,6} {\draw (\x,0.12)--(\x,-0.12);}
        \draw[thick,green!60!black] (3,0)--(3,1.8);
        \fill[green!60!black] (3,1.8) circle (2.1pt);
        \node[green!60!black] at (3,-0.35) {$K_0+1$};
        \node[green!60!black] at (1.6,2.05) {$u_{K_0+1}$};
        \foreach \x in {1,2,4,5,6} {\fill[green!60!black] (\x,0) circle (2pt);}
    \end{scope}

    % asse 3
    \begin{scope}[yshift=-1.4cm]
        \draw[->] (0,0) -- (7.2,0) node[right] {$k$};
        \draw[->] (0,0) -- (0,2.4) node[above] {$U$};
        \foreach \x in {1,2,3,4,5,6} {\draw (\x,0.12)--(\x,-0.12);}
        \draw[thick,violet] (4,0)--(4,1.4);
        \fill[violet] (4,1.4) circle (2.1pt);
        \node[violet] at (4,-0.35) {$K_0+2$};
        \node[violet] at (2.5,1.65) {$u_{K_0+2}$};
        \foreach \x in {1,2,3,5,6} {\fill[violet] (\x,0) circle (2pt);}
    \end{scope}
\end{tikzpicture}
\caption{Esempi di ingressi elementari: ciascuno è nullo ovunque tranne che in un solo istante.}
\end{figure}

\section{Dalla decomposizione dell'ingresso alla risposta forzata nel tempo discreto}

La risposta forzata è lineare rispetto all'ingresso. Applicando la decomposizione \eqref{eq:L3_input_decomposition} si ottiene
\[
\phi_f(K,K_0,u)=\sum_{j=K_0}^{K-1}\phi_f(K,K_0,u_j).
\]

Ora osserviamo che il segnale $u_j$ è nullo fino all'istante $j-1$ e si attiva solo all'istante $j$.  
Per causalità e composizione, il suo contributo allo stato a tempo $K$ dipende solo da ciò che accade a partire da $j$, quindi
\[
\phi_f(K,K_0,u_j)=\phi_f(K,j,u_j).
\]

Poiché $u_j$ è determinato dal solo vettore $u(j)\in U$, esiste una matrice
\[
H(K,j)\in\mathbb{R}^{n\times m}
\]
tale che
\[
\phi_f(K,j,u_j)=H(K,j)\,u(j).
\]

Sostituendo:
\[
\phi_f(K,K_0,u)=\sum_{j=K_0}^{K-1}H(K,j)\,u(j).
\]

Infine, sommando risposta libera e risposta forzata:
\[
x(K)=\Phi(K,K_0)x(K_0)+\sum_{j=K_0}^{K-1}H(K,j)\,u(j).
\]

\paragraph{Significato della matrice $H(K,j)$.}
La matrice $H(K,j)$ descrive come un ingresso applicato all'istante $j$ contribuisce allo stato osservato all'istante finale $K$.

\section{Rappresentazione esplicita discreta e passaggio alla forma ricorsiva}

La formula
\[
x(K)=\Phi(K,K_0)x(K_0)+\sum_{j=K_0}^{K-1}H(K,j)\,u(j)
\]
è una \emph{rappresentazione esplicita} del sistema a tempo discreto.

Se ora scegliamo
\[
K=j+1,
\]
otteniamo
\[
x(j+1)=\Phi(j+1,j)x(j)+H(j+1,j)u(j).
\]

Definendo
\[
A(j):=\Phi(j+1,j),
\qquad
B(j):=H(j+1,j),
\]
si ricava la forma ricorsiva standard
\begin{equation}
x(j+1)=A(j)x(j)+B(j)u(j).
\label{eq:L3_TD_implicit}
\end{equation}

L'uscita rimane
\begin{equation}
y(j)=C(j)x(j)+D(j)u(j).
\label{eq:L3_TD_output_local}
\end{equation}

\paragraph{Conclusione.}
Nel tempo discreto:
\begin{itemize}
    \item la rappresentazione con $\Phi$ e $H$ è la forma \emph{esplicita};
    \item la rappresentazione con $A(j)$ e $B(j)$ è la forma \emph{implicita} o ricorsiva.
\end{itemize}

\section{Proprietà di consistenza e composizione di \texorpdfstring{$\Phi$}{Phi}}

Le proprietà della funzione di transizione $\phi$ si traducono in proprietà matriciali di $\Phi$ e $H$.

\subsection{Consistenza}

Per definizione di consistenza:
\[
\phi(t_0,t_0,x(t_0),u)=x(t_0).
\]

Nel caso lineare e libero:
\[
x(t_0)=\Phi(t_0,t_0)x(t_0)\qquad \forall x(t_0).
\]

Poiché questa relazione deve valere per ogni stato iniziale, segue che
\begin{equation}
\Phi(t_0,t_0)=I_n.
\label{eq:L3_consistency_Phi}
\end{equation}

La stessa proprietà vale sia nel tempo continuo sia nel tempo discreto:
\[
\Phi(t,t)=I_n,
\qquad
\Phi(K,K)=I_n.
\]

\subsection{Composizione per \texorpdfstring{$\Phi$}{Phi}}

Dalla proprietà generale di composizione della dinamica:
\[
\phi(t_2,t_0,x_0,u)=\phi\bigl(t_2,t_1,\phi(t_1,t_0,x_0,u),u\bigr),
\qquad t_0\le t_1\le t_2,
\]
e scegliendo ingresso nullo $u=0$, si ottiene
\[
\Phi(t_2,t_0)x_0=\Phi(t_2,t_1)\Phi(t_1,t_0)x_0.
\]

Essendo vero per ogni $x_0$, segue:
\begin{equation}
\Phi(t_2,t_0)=\Phi(t_2,t_1)\Phi(t_1,t_0),
\qquad t_0\le t_1\le t_2.
\label{eq:L3_composition_Phi_TC}
\end{equation}

Nel discreto:
\begin{equation}
\Phi(K_2,K_0)=\Phi(K_2,K_1)\Phi(K_1,K_0),
\qquad K_0\le K_1\le K_2.
\label{eq:L3_composition_Phi_TD}
\end{equation}

\paragraph{Interpretazione.}
La matrice di transizione da $t_0$ a $t_2$ si ottiene componendo la transizione da $t_0$ a $t_1$ con quella da $t_1$ a $t_2$.

\section{Proprietà di composizione di \texorpdfstring{$H$}{H}}

Usiamo ora la rappresentazione esplicita discreta in due modi diversi.

Da un lato:
\[
x(K_2)=\Phi(K_2,K_0)x(K_0)+\sum_{j=K_0}^{K_2-1}H(K_2,j)u(j).
\]

Dall'altro, applicando prima la dinamica da $K_0$ a $K_1$ e poi da $K_1$ a $K_2$:
\[
x(K_2)=\Phi(K_2,K_1)x(K_1)+\sum_{j=K_1}^{K_2-1}H(K_2,j)u(j).
\]

Sostituiamo nella seconda formula
\[
x(K_1)=\Phi(K_1,K_0)x(K_0)+\sum_{j=K_0}^{K_1-1}H(K_1,j)u(j),
\]
ottenendo
\[
x(K_2)=
\Phi(K_2,K_1)\Phi(K_1,K_0)x(K_0)
+
\sum_{j=K_0}^{K_1-1}\Phi(K_2,K_1)H(K_1,j)u(j)
+
\sum_{j=K_1}^{K_2-1}H(K_2,j)u(j).
\]

Confrontando con la formula diretta e usando l'unicità dei coefficienti, si ottengono due risultati:

\begin{enumerate}
    \item la proprietà di composizione di $\Phi$ già vista;
    \item la proprietà di composizione di $H$:
    \begin{equation}
    \Phi(K_2,K_1)H(K_1,j)=H(K_2,j),
    \qquad K_0\le j<K_1\le K_2.
    \label{eq:L3_composition_H_TD}
    \end{equation}
\end{enumerate}

Nel tempo continuo l'analogo è
\begin{equation}
\Phi(t_2,t_1)H(t_1,\tau)=H(t_2,\tau),
\qquad t_0\le \tau \le t_1\le t_2.
\label{eq:L3_composition_H_TC}
\end{equation}

\paragraph{Interpretazione.}
Per propagare fino al tempo finale l'effetto di un ingresso applicato in un istante precedente, posso:
\begin{itemize}
    \item calcolarne prima l'effetto fino a un tempo intermedio;
    \item poi trasportare quell'effetto dal tempo intermedio a quello finale con $\Phi$.
\end{itemize}

\section{Riassunto strutturale delle rappresentazioni esplicite}

Le proprietà appena ricavate permettono di scrivere in modo ordinato le rappresentazioni esplicite dei sistemi lineari.

\subsection{Tempo discreto}

\begin{equation}
\boxed{
\begin{aligned}
x(K) &= \Phi(K,K_0)x(K_0)+\sum_{j=K_0}^{K-1}H(K,j)u(j),\\
y(K) &= C(K)x(K)+D(K)u(K).
\end{aligned}
}
\label{eq:L3_explicit_discrete_boxed}
\end{equation}

\subsection{Tempo continuo}

\begin{equation}
\boxed{
\begin{aligned}
x(t) &= \Phi(t,t_0)x(t_0)+\int_{t_0}^{t}H(t,\tau)u(\tau)\,d\tau,\\
y(t) &= C(t)x(t)+D(t)u(t).
\end{aligned}
}
\label{eq:L3_explicit_continuous_boxed}
\end{equation}

\section{Come si passa nel continuo alla forma implicita}

Negli appunti vengono poi ricavate due relazioni fondamentali, che permettono di passare dalla rappresentazione esplicita continua a quella implicita.

\subsection{Primo lemma: derivata della matrice di transizione}

Usiamo la proprietà di composizione:
\[
\Phi(t+\Delta t,t_0)=\Phi(t+\Delta t,t)\,\Phi(t,t_0).
\]

Sottraendo $\Phi(t,t_0)$ da entrambi i membri:
\[
\Phi(t+\Delta t,t_0)-\Phi(t,t_0)
=
\bigl[\Phi(t+\Delta t,t)-I_n\bigr]\Phi(t,t_0).
\]

Dividendo per $\Delta t$ e passando al limite per $\Delta t\to 0$:
\[
\frac{\partial \Phi(t,t_0)}{\partial t}
=
\left[
\frac{\partial \Phi(\zeta,t)}{\partial \zeta}
\right]_{\zeta=t}
\Phi(t,t_0).
\]

Definiamo allora
\begin{equation}
A(t):=
\left[
\frac{\partial \Phi(\zeta,t)}{\partial \zeta}
\right]_{\zeta=t}.
\label{eq:L3_def_A}
\end{equation}

Quindi il primo lemma diventa
\begin{equation}
\frac{\partial \Phi(t,t_0)}{\partial t}=A(t)\Phi(t,t_0).
\label{eq:L3_lemma1}
\end{equation}

\subsection{Secondo lemma: derivata del nucleo forzato}

Usando la proprietà di composizione di $H$,
\[
H(t+\Delta t,\tau)=\Phi(t+\Delta t,t)\,H(t,\tau),
\]
si ottiene in modo analogo:
\[
\frac{H(t+\Delta t,\tau)-H(t,\tau)}{\Delta t}
=
\frac{\Phi(t+\Delta t,t)-I_n}{\Delta t}\,H(t,\tau).
\]

Passando al limite:
\begin{equation}
\frac{\partial H(t,\tau)}{\partial t}=A(t)H(t,\tau).
\label{eq:L3_lemma2}
\end{equation}

\section{Derivazione della forma implicita nel tempo continuo}

Partiamo dalla rappresentazione esplicita
\[
x(t)=\Phi(t,t_0)x(t_0)+\int_{t_0}^{t}H(t,\tau)u(\tau)\,d\tau.
\]

Deriviamo rispetto a $t$.

\subsection{Derivata del primo termine}

Usando il lemma \eqref{eq:L3_lemma1}:
\[
\frac{d}{dt}\Bigl(\Phi(t,t_0)x(t_0)\Bigr)=A(t)\Phi(t,t_0)x(t_0).
\]

\subsection{Derivata del termine integrale}

Applichiamo la regola di Leibniz:
\[
\frac{d}{dt}\left(
\int_{t_0}^{t}H(t,\tau)u(\tau)\,d\tau
\right)
=
H(t,t)u(t)+\int_{t_0}^{t}\frac{\partial H(t,\tau)}{\partial t}u(\tau)\,d\tau.
\]

Usando il lemma \eqref{eq:L3_lemma2}:
\[
=
H(t,t)u(t)+\int_{t_0}^{t}A(t)H(t,\tau)u(\tau)\,d\tau.
\]

Poiché $A(t)$ non dipende dalla variabile di integrazione $\tau$, si può portare fuori dall'integrale:
\[
=
H(t,t)u(t)+A(t)\int_{t_0}^{t}H(t,\tau)u(\tau)\,d\tau.
\]

\subsection{Ricomposizione finale}

Sommiamo i due contributi:
\[
\dot x(t)
=
A(t)\Phi(t,t_0)x(t_0)
+
A(t)\int_{t_0}^{t}H(t,\tau)u(\tau)\,d\tau
+
H(t,t)u(t).
\]

Raccogliendo $A(t)$:
\[
\dot x(t)
=
A(t)\left[
\Phi(t,t_0)x(t_0)+\int_{t_0}^{t}H(t,\tau)u(\tau)\,d\tau
\right]
+
H(t,t)u(t).
\]

Ma la quantità tra parentesi quadre è proprio $x(t)$, quindi
\[
\dot x(t)=A(t)x(t)+H(t,t)u(t).
\]

Definendo
\begin{equation}
B(t):=H(t,t),
\label{eq:L3_def_B}
\end{equation}
si ottiene finalmente la forma implicita del sistema lineare continuo:
\begin{equation}
\boxed{
\dot x(t)=A(t)x(t)+B(t)u(t).
}
\label{eq:L3_TC_implicit}
\end{equation}

L'uscita è
\begin{equation}
\boxed{
y(t)=C(t)x(t)+D(t)u(t).
}
\label{eq:L3_TC_output_implicit}
\end{equation}

\section{Relazione tra forma esplicita e forma implicita}

A questo punto il quadro è completo.

\subsection{Tempo discreto}

\paragraph{Forma esplicita}
\[
x(K)=\Phi(K,K_0)x(K_0)+\sum_{j=K_0}^{K-1}H(K,j)u(j).
\]

\paragraph{Forma implicita}
\[
x(k+1)=A(k)x(k)+B(k)u(k),
\]
con
\[
A(k)=\Phi(k+1,k),\qquad B(k)=H(k+1,k).
\]

\subsection{Tempo continuo}

\paragraph{Forma esplicita}
\[
x(t)=\Phi(t,t_0)x(t_0)+\int_{t_0}^{t}H(t,\tau)u(\tau)\,d\tau.
\]

\paragraph{Forma implicita}
\[
\dot x(t)=A(t)x(t)+B(t)u(t),
\]
con
\[
A(t)=\left[\frac{\partial \Phi(\zeta,t)}{\partial \zeta}\right]_{\zeta=t},
\qquad
B(t)=H(t,t).
\]

\paragraph{Messaggio concettuale della lezione.}
Le matrici $\Phi$ e $H$ descrivono l'evoluzione \emph{globale} del sistema:
\begin{itemize}
    \item $\Phi$ trasporta nel tempo lo stato iniziale;
    \item $H$ pesa i contributi dell'ingresso applicato ai vari istanti.
\end{itemize}

Le matrici $A$ e $B$, invece, descrivono la legge \emph{locale} di evoluzione:
\begin{itemize}
    \item $A$ governa la dinamica interna;
    \item $B$ governa il modo in cui l'ingresso agisce istante per istante sullo stato.
\end{itemize}

\section{Sintesi finale della lezione}

In questa lezione sono stati introdotti e chiariti i seguenti punti:
\begin{enumerate}
    \item nei sistemi lineari l'uscita si scrive come
    \[
    y=Cx+Du;
    \]
    \item la risposta libera dello stato è rappresentata dalla matrice di transizione $\Phi$;
    \item la risposta forzata è descritta da un nucleo $H$;
    \item nel tempo discreto l'evoluzione forzata diventa una somma finita grazie alla decomposizione dell'ingresso in segnali elementari;
    \item $\Phi$ e $H$ soddisfano proprietà di consistenza e composizione;
    \item nel continuo, derivando la forma esplicita, si ottiene la forma implicita
    \[
    \dot x=A(t)x+B(t)u(t).
    \]
\end{enumerate}

Questa struttura è fondamentale perché nei capitoli successivi permetterà di studiare:
\begin{itemize}
    \item esponenziale di matrice;
    \item evoluzione libera e forzata;
    \item stabilità;
    \item raggiungibilità e osservabilità dei sistemi lineari.
\end{itemize}